\color{black}
\section{Case study: pancreatic cancer}\label{section:PDAC}

Pancreatic ductal adenocarcinoma (PDAC) is among the most aggressive and deadly cancers, characterised by rapid progression, late-stage detection, and  frequent liver metastasis \citep{Pereira2019Metastasis}. 
The long-term outlook for patients remains poor, with median survival below six months and a five-year survival rate under 5\% \citep{kamisawa2016pancreatic, li2019angiogenesis}. 
PDAC exhibits substantial heterogeneity across patients, as highlighted by The Cancer Genome Atlas (TCGA) project \citep{Sidaway2017TCGA}.
Unlike many other malignancies, PDAC typically remains asymptomatic until it reaches an advanced stage. 
Surgical resection is therefore often the only curative option, typically followed by adjuvant radiation therapy. 

We study the causal effect of adjuvant radiotherapy on survival while adjusting for genetic and clinical factors that may confound or modify this effect.
We analyse data from the TCGA-PAAD cohort \citep{Sidaway2017TCGA}.
The dataset contains gene expression and clinical information for 181 PDAC tumours.
Data are curated via the \texttt{pdacR} package \citep{TorreHealy2023pdacR}.
We excluded 51 patients with missing values in either survival time, censoring status, or treatment assignment.
This leaves $ 130$ patients.
The outcome is overall survival, subject to right-censoring. 
The observed censoring rate  is 47\%.

The analysis includes several clinical covariates (see Supplementary Materials~\ref{appendix:clinical_covariates}, \citealp{supplement}) and gene expression profiles. 
A number of well-established genetic alterations drive PDAC tumorigenesis \citep{Stefanoudakis2024}.
We a priori include a subset of genes previously linked to PDAC in the literature (see Supplementary Materials~\ref{appendix:clinical_covariates}).
Additionally, we incorporate the 3000 genes with the highest median absolute deviation (MAD) across patients to capture the most variable and remove low-variability genes.
This approach emphasises genes with substantial variation, which are more likely to carry informative signal for downstream modelling and reduces noise from uninformative or uniformly expressed genes \citep{Hackstadt2009, Zhang2018}.


Median follow-up time is 23.3 months.
This is computed using the reverse Kaplan--Meier method \citep{Schemper1996}.
The estimated 2-year overall survival probability is 39.2\%.  
Figure~\ref{fig:PDAC_KaplanMeier} shows Kaplan--Meier survival curves stratified by treatment arm.

\input PDAC_KM 

We fit a Causal Horseshoe Forest to the data with aggressive default settings in Section~\ref{default_hyperparameters}.
We performed stratified $10 \times 5$-fold cross-validation to validate the choice of hyperparameter $k$. 
Cross-validation showed no difference in terms of predictive performance across different choices of $k$ (see Supplementary Materials~\ref{appendix:clinical_covariates}).
We assessed convergence using trace plots of $\sigma^2$ (see Supplementary Materials~\ref{appendix:clinical_covariates}).
After a burn-in of 5000 iterations, we retained 5000 posterior samples for inference. 
The fitted model achieved a concordance index of 0.80, indicating good performance.

\input PDAC_ATE_CATEs 


Figure~\ref{fig:posterior_ATE} shows the posterior distribution of the  average treatment effect (ATE, see Supplementary Materials~\ref{appendix:clinical_covariates}), with the 95\% credible interval marked by vertical dashed lines.
The posterior mean ATE is approximately 0.75, with a 95\% credible interval ranging from 0.40 to 1.11. 
This distribution is concentrated above zero, providing strong evidence for a positive average treatment effect.
These results suggest that adjuvant radiation therapy improves overall survival on average in PDAC patients.

Figure~\ref{fig:cate_plot} shows the posterior CATEs for each individual.
Each horizontal line represents a 95\% credible interval for the conditional average treatment effect.
Most intervals are wide and include zero.
This reflects substantial uncertainty and limited evidence for strong individualised effects.
A small subset of patients has intervals entirely above zero.
This suggests evidence---if any---for a positive treatment effect.
The wide intervals for many patients likely reflect insufficient information to confidently distinguish their individualised treatment effects for individuals with similar characteristics.
Nevertheless, the consistent rightward shift of the posterior means supports an overall beneficial and homogeneous effect of the treatment.



We summarise treatment effect heterogeneity using posterior probabilities of differential treatment effect and benefit \citep{henderson2020individualized}. 
For each individual, we compute 
$D_i = \p(\tau(X_i) \ge \bar\tau \mid \mathcal D)$, 
where $\bar\tau$ denotes the model-based average treatment effect and $\mathcal{D}$ the observed data. 
We transform this to 
$D_i^\ast = \max\{1 - 2D_i, 2D_i - 1\} \in [0,1]$, 
which measures posterior evidence for deviation from the population average. 
Values near $0$ indicate little evidence of heterogeneity, whereas values near $1$ indicate strong evidence. 
Population-level heterogeneity is summarised by reporting the proportion of individuals for whom $D_i^\ast$ exceeds pre-specified thresholds. 
We additionally report posterior probabilities of treatment benefit, 
$\p(\tau(X_i) > 0 \mid \mathcal D)$.

Table~\ref{tab:pdac_heterogeneity_summary} shows little evidence of substantial treatment effect heterogeneity in this cohort. 
No individuals exhibit mild or strong posterior evidence of a differential effect relative to the average treatment effect. 
At the same time, posterior probabilities of benefit are uniformly high: approximately $96.9\%$ of patients satisfy $\p(\tau(X_i) > 0 \mid \mathcal D) > 0.95$, with the remainder falling in the $(0.75, 0.95]$ range. 
These findings are consistent with earlier results: a positive average treatment effect with limited evidence for meaningful between-patient variation beyond posterior uncertainty.



\begin{table}[b!]
\centering
\begin{tabular}{lc}
\toprule
Summary measure & Percentage (\%) \\
\midrule
$D_i^\ast > 0.95$ & 0.0 \\
$D_i^\ast > 0.80$ & 0.0 \\
$\p(\tau(x_i) > 0 \mid \mathcal D) \in (0.95, 1]$ & 96.9 \\
$\p(\tau(x_i) > 0 \mid \mathcal D) \in (0.75, 0.95]$ & 3.1 \\
$\p(\tau(x_i) > 0 \mid \mathcal D) \in (0.25, 0.75]$ & 0.0 \\
$\p(\tau(x_i) > 0 \mid \mathcal D) \in [0, 0.25]$ & 0.0 \\
\bottomrule
\end{tabular}
    \caption[Posterior summaries of treatment effect heterogeneity]{Posterior summaries of treatment effect heterogeneity and evidence of benefit in the PDAC analysis. The first two rows report the proportion of patients with posterior evidence of differential treatment effect relative to the average, measured by $D_i^\ast$. \citet{henderson2020individualized} suggest $D_i^\ast>0.95$ as strong evidence of differential treatment effect and $D_i^\ast>0.80$ as mild evidence.}
  \label{tab:pdac_heterogeneity_summary}
\end{table}

Stable Unit Treatment Value Assumption (SUTVA) is not fully met, as adjuvant radiation therapy is not standardised across patients.
Treatment regimens may vary in composition, dosage, and duration.
As a result, our analysis estimates the effect of receiving any form of radiation therapy, rather than a single, uniform intervention.
We assume sufficient overlap in a lower-dimensional confounding subspace. 
Strict overlap is unlikely in the full covariate space given the data dimensions \citep{DAmour2021Overlap}.
Positivity is fundamentally unverifiable in the $p \gg n$ regime considered here. 
Any estimated propensity score is itself a heavily regularised, high-dimensional object. 
Standard diagnostics such as propensity score histograms reflect this shrinkage behaviour as much as structural overlap. 
We therefore treat overlap as a working assumption that the data cannot directly verify.


The high mortality rate of pancreatic cancer introduces bias.
Some patients may die before becoming eligible for or completing adjuvant radiation therapy.
Thus, patients in the treatment group must have survived long enough to initiate and complete therapy.
This creates immortal time \citep{Levesque2010ImmortalTimeBias}, a period during which the event cannot occur by definition, potentially biasing results in favour of the treatment group.
We provide a preliminary landmark analysis \citep{Anderson1983TumorResponse, Putter2013Landmarking} in Supplementary Materials~\ref{appendix:landmark}.
A more thorough analysis would require detailed data on treatment timing, which are not available in this study. 
Future studies with detailed treatment timing data could help to address this bias.


